%% Predictive Coding: What Is Exact, What Is Not, and What It Costs
%%
%% Machine-checked theory of predictive-coding dynamics (Lean 4, Mettapedia
%% library), with the experimental verdict it explains.
%% Companion empirical report: AlienCoding.tex.
%%
%% Build: pdflatex PC.tex

\documentclass[11pt]{article}
\usepackage[T1]{fontenc}
\usepackage[utf8]{inputenc}
\usepackage{lmodern}
\usepackage{geometry}
\geometry{margin=1in}
\usepackage{hyperref}
\usepackage{booktabs}
\usepackage{amsmath}
\usepackage{amssymb}
\usepackage{amsthm}
\usepackage{float}
\usepackage{longtable}
\usepackage{array}
\usepackage{fvextra}
\usepackage{newunicodechar}
\usepackage{needspace}
\setlength{\emergencystretch}{3em}

\newunicodechar{ℝ}{\ensuremath{\mathbb{R}}}
\newunicodechar{ℕ}{\ensuremath{\mathbb{N}}}
\newunicodechar{ℂ}{\ensuremath{\mathbb{C}}}
\newunicodechar{∀}{\ensuremath{\forall}}
\newunicodechar{∃}{\ensuremath{\exists}}
\newunicodechar{∧}{\ensuremath{\land}}
\newunicodechar{∨}{\ensuremath{\lor}}
\newunicodechar{↔}{\ensuremath{\leftrightarrow}}
\newunicodechar{≤}{\ensuremath{\leq}}
\newunicodechar{≠}{\ensuremath{\neq}}
\newunicodechar{∑}{\ensuremath{\sum}}
\newunicodechar{∫}{\ensuremath{\int}}
\newunicodechar{∂}{\ensuremath{\partial}}
\newunicodechar{𝓝}{\ensuremath{\mathcal{N}}}
\newunicodechar{ᶠ}{\ensuremath{^{\mathrm f}}}
\newunicodechar{⁻}{\ensuremath{^{-}}}
\newunicodechar{ₗ}{\ensuremath{_{l}}}

\newtheorem{theorem}{Theorem}[section]
\newtheorem{corollary}[theorem]{Corollary}
\newcommand{\leanname}[1]{{\scriptsize\ttfamily
  \def\_{\textunderscore\allowbreak}#1}}
\newcommand{\leancommit}{647168959fc57cfe97526d9525d33e63681b8e2f}
\newcommand{\leansource}[3]{%
  \href{https://github.com/zariuq/MeTTapedia/blob/\leancommit/#1\#L#2}{\leanname{#3}}}
\newcommand{\leansignaturelink}[2]{%
  \par\Needspace{0.3\textheight}\smallskip\noindent
  \href{https://github.com/zariuq/MeTTapedia/blob/\leancommit/#1\#L#2}{%
    \scriptsize\itshape Lean 4 signature and complete proof}\par\nobreak}
\DefineVerbatimEnvironment{leansignature}{Verbatim}{%
  fontsize=\scriptsize,breaklines=true,breakanywhere=true,
  frame=single,framesep=1.5mm,xleftmargin=0.5em,xrightmargin=0.5em}

\title{Predictive Coding: What Is Exact,\\ What Is Not, and What It Costs}
\author{Oru\v{z}i\thanks{Oru\v{z}i names the AI research collaborators who carried out
this work: Anthropic Claude (Opus~4 and Fable~5) and OpenAI Codex (GPT-5.6-Sol) ---
formalization, implementation, experiments, and drafting. Research direction and
oversight: Zar Goertzel.}}
\date{July 2026 --- working note}

\begin{document}
\maketitle

\begin{abstract}
We prove what predictive coding (PC) does and does not compute, and test the
consequences. Settling any finite linear-Gaussian residual model is exactly conditional
Bayesian inference --- the unique energy equilibrium \emph{is} the posterior mean, at
arbitrary width, depth, and correlated precision. The equilibrium error recursion is
exactly backpropagation's, but the equilibrium weight gradient is not, and we give the
counterexample. Exact backprop-equivalent scheduling is characterized completely: a
schedule is exact iff the contributions it misses sum to zero, so when contributions
cannot cancel it reduces to a structural criterion --- every node's outgoing edges land
on a single rank --- strictly weaker than the levelizability one expects. Error-based PC
is an exact nonlinear change of coordinates from state-based PC, with equal energy and
local parameter updates; its first-arrival signal carries exactly the claimed
$\lambda^{L-i}$ factor. Composing that factor with the local update and loss laws
proves a depth-credit no-go: inference rates equalize every layer only at the unit-rate
backprop boundary, while polynomial Depth-$\mu$P scaling cannot cancel exponential
attenuation. Settling on a depth-$n$ path contracts at exactly
$\cos(\pi/n)$.  For a genuinely joined backbone--adapter path, restricting the
coupled operator to a frozen backbone yields exactly the adapter-only first-mode rate;
leaving the backbone active re-imports its depth bound.  A preconditioned PC step
multiplies the loss by $(P-1)^2$: descent holds exactly on $0<P<2$, only $P=1$ reaches
zero in one step, and asymptotically all that is required is $\rho(I-P)<1$ --- which a
nonnormal Jordan family shows can coexist with arbitrarily large finite transients.
The unrestricted conjecture that every saddle of equilibrated PC energy is strict is
false: we give a critical cubic-order saddle with zero curvature in every direction.
A separately trained centered-sigmoid PC model has exact and robust entropy turnover,
while also proving contraction alone insufficient for reachability. Quadratic PC energy
cannot orient even a two-node edge from any observational dataset; ordinary clamping is
conditioning in both directions. A uniform intervention --- clamp a node and delete its
own residual term --- separates the directions exactly when their covariance is nonzero.
On any excited finite two-node intervention design, after both orientations minimize
over their own structural coefficient, the true orientation has strictly lower energy
if and only if that covariance is nonzero.
Finally, two priors with \emph{equal}
cross-entropy can have strictly ordered expected search yields: fit does not imply
guidance. All statements are machine-checked in Lean~4. Experimentally, PC never
overtook backpropagation in any regime we measured: every from-scratch arm failed its
registered bar, and the one ignition that reproduced across seeds had to start from a
backprop-equivalent anchor.
\end{abstract}

\section{Introduction}

Predictive coding networks train by relaxing latent activities against local prediction
errors, then applying purely local weight updates. The hope is a biologically plausible
--- and practically competitive --- alternative to backpropagation, and the literature
supplies equivalence results in the settling limit together with a growing engineering
effort to make deep PC trainable at all
\cite{whittington2017,song2020,millidge2020,salvatori2022,goemaere2025}.

Those equivalence results are stated informally, and they are sloppiest exactly where
the hypotheses matter: \emph{which} quantity matches backpropagation, \emph{under what
schedule}, \emph{at what depth}, and \emph{at what cost}. This note answers those four
questions with theorems, in models small enough that every hypothesis is explicit and
every boundary carries a counterexample: linear chains, finite DAGs, scalar paths,
finite categorical priors. Each is machine-checked in Lean~4 and axiom-clean, using only
\leanname{propext}, \leanname{Classical.choice} and \leanname{Quot.sound}. Selected
crown results display their source-exact Lean theorem headers in place, with proof
bodies omitted; each caption links to the complete declaration and proof in an
immutable source snapshot. Table~\ref{tab:index} links every supporting statement and
boundary fixture in the same snapshot.

Two of the answers surprised us. The natural conjecture about which graphs admit exact
scheduling is false in an instructive direction (Section~\ref{sec:schedule}), and the
question every PC-versus-BP comparison reports on --- training loss --- provably fails
to determine the thing those comparisons care about (Section~\ref{sec:yield}).

Section~\ref{sec:experiments} reports the experiments the theory explains. Every arm
fixed its acceptance criteria before running; failed arms are reported as failures with
their measured failure modes. The companion report \emph{AlienCoding} carries the full
empirical account.

\section{The model}

A predictive-coding network attaches to each node $i$ of a computation graph a latent
state $z_i$, a prediction $f_i(z_{\mathrm{pa}(i)};w)$ from its parents, and a precision
$\lambda_i > 0$. The energy is
\[
E(z,w) \;=\; \sum_i \lambda_i\,\bigl(z_i - f_i(z_{\mathrm{pa}(i)};w)\bigr)^2 .
\]
\emph{Settling} clamps the inputs (and, in supervised mode, the outputs) and descends
$E$ in $z$; the \emph{weight update} is the local gradient $-\partial E/\partial w$
evaluated at the settled state. Backpropagation instead differentiates the supervised
loss through the forward activations $a$. Every equivalence claim in the field is,
underneath, a claim about the relationship between the settled state $z^\star$ and the
forward state $a$.

Our running instance is the scalar chain of depth $d$: latents $z_0,\dots,z_d$, links
$i < d$ with gain $w_i$ and precision $\lambda_i$, residuals and energy
\[
r_i \;=\; z_{i+1} - w_i z_i, \qquad E(z,w) \;=\; \sum_{i<d} \lambda_i r_i^2 ,
\]
clamped at $z_0 = x$ and $z_d = y$. We carry no factor of $\tfrac12$, so the weight
gradient is $\partial E/\partial w_i = -2\,e_i\,z_i$ where
\[
e_i \;:=\; \lambda_i r_i
\]
is the precision-weighted error at link $i$. The quantity $e_i$ is the one that will
turn out to be exact.

\section{Settling is Bayesian inference}

Read the chain as a linear-Gaussian generative model: each $z_{i+1}$ is $w_i z_i$ plus
noise of precision $\lambda_i$. Clamping $z_0$ and $z_d$ conditions on the two ends.

The right level for this is not the chain. Take any finite \emph{affine residual model}:
latents $u$, residual $r(u) = Au + c$ with $A$ injective, and an arbitrary
positive-definite residual precision $\Lambda$ --- so residuals may be correlated, and
$A$ need not be square. The energy is $r(u)^\top \Lambda\, r(u)$.

\begin{theorem}[Settling computes the posterior mean]\label{thm:bayes}
For any such model the posterior precision $Q = A^\top \Lambda A$ is positive definite,
and the following coincide: $u$ is an energy equilibrium; $u$ is the unique energy
minimizer; and $u$ is the conditional multivariate-Gaussian posterior mean. Any clamped
state other than the posterior mean is provably not an equilibrium.
\end{theorem}

\leansignaturelink{lean/mettapedia/Mettapedia/MachineLearning/NeuralNetworks/%
PredictiveCoding/LinearGaussianOperator.lean}{254}
\begin{leansignature}
theorem LinearGaussianOperatorModel.equilibrium_iff_eq_conditionalPosteriorMean
    (model : LinearGaussianOperatorModel Latent Residual)
    (u : LinearGaussianOperatorSpace Latent) :
    model.Equilibrium u ↔ u = ∫ v, v ∂model.posterior
\end{leansignature}

So ``settling performs Bayesian inference'' is a theorem, not a slogan: the relaxed
latents \emph{are} the posterior mean given the evidence. Scalar chains and
vector-valued chains --- arbitrary width, arbitrary depth, arbitrary positive-definite
(hence correlated) residual precision --- are corollaries, obtained by exhibiting the
chain's residual map as such an $A$ and checking injectivity; no chain-specific
completion-of-squares argument is repeated. At depth two the scalar case recovers the
Kalman update and the two-source Gaussian fusion rule, which is where the PC
literature's Bayesian intuition comes from.

This is also the cleanest way to see the trouble in the next section. The posterior mean
given a clamped output is precisely \emph{not} the forward pass --- conditioning on the
answer moves the latents. That is the point of PC, and it is also why its weight gradient
is not backpropagation's.

\section{What matches backpropagation, and what does not}

\subsection{The error recursion is exact}

\begin{theorem}[Equilibrium errors obey backprop's recursion]\label{thm:errrec}
At any clamped equilibrium of the scalar chain, for every $i$ with $i+1 < d$,
\[
e_i \;=\; w_{i+1}\, e_{i+1}.
\]
\end{theorem}

This is backpropagation's backward recursion, exactly, with no limit and no
approximation: errors travel from the output end to the input end multiplied by the
intervening gains. Weight tying preserves it, which is the credit-assignment certificate
for recurrent PC.

The same statement survives the move to the graphs real implementations use. Hash-consed
evaluation DAGs share subexpressions, and the design question is whether a shared node
gets one latent or one per use.

\begin{theorem}[Shared-node aggregation]\label{thm:dagrec}
Let $G$ be a finite DAG with one latent per node, edge gains $w_e$, node precisions
$\lambda_v$, and errors $e_v = \lambda_v r_v$. At any equilibrium, every unclamped node
$v$ satisfies
\[
e_v \;=\; \sum_{e\,:\,\mathrm{src}(e) = v} w_e\, e_{\mathrm{tgt}(e)} .
\]
\end{theorem}

That is reverse-mode accumulation: a shared node's settled error is the gain-weighted
sum of the errors of everything it feeds, which is what differentiation does when it
sums gradients over the unfolded tree. The negative fixture is the one that matters to
implementors: on a graph where node~$0$ feeds node~$1$ through two edges, the correct
settled error at node~$0$ is $2$, whereas an implementation retaining a single parent
slot records~$1$ --- half the true value. One latent per shared node, all slots summed,
is forced. Our C implementation of PC for the tree network (Section~\ref{sec:tree}) is
built this way, and its numerical oracles test these shapes.

\subsection{The gradient identity is false}

The frequently repeated stronger claim is that the equilibrium \emph{weight gradient}
equals the backprop gradient. It does not, and linearity does not save it.

\begin{theorem}[Counterexample]\label{thm:gradne}
Take the depth-2 chain with $w_0 = w_1 = 1$, $\lambda_0 = \lambda_1 = 1$, input $x=1$
and clamped target $y=2$. The forward state is $a = (1,1,1)$ and the clamped equilibrium
is $z^\star = (1,\tfrac32,2)$. Then
\[
\frac{\partial E}{\partial w_0}\Big|_{z^\star} = -1
\qquad\text{while}\qquad
\text{backprop's gradient for } w_0 = -2 .
\]
The root cause is $z^\star_1 = \tfrac32 \neq 1 = a_1$.
\end{theorem}

\leansignaturelink{lean/mettapedia/Mettapedia/MachineLearning/NeuralNetworks/%
PredictiveCoding/BackpropExactness.lean}{101}
\begin{leansignature}
theorem equilibriumGradient_ne_backpropGradient_example :
    pcGainPartial unitDepth2Links depth2EquilibriumState 0 ≠
      scalarSquaredErrorBackpropPartial 1 1 1
\end{leansignature}

Both gradients use the correct error; they evaluate it against different source
activations. Backprop multiplies by the forward activation $a_1$; PC multiplies by the
relaxed latent $z^\star_1$, which conditioning has already moved toward the answer. On
this fixture the two differ by a factor of two, and nothing about the linear case
prevents it --- by Theorem~\ref{thm:bayes} the displacement $z^\star \neq a$ is exactly
the Bayesian content of settling.

The consequence for the field is a naming discipline. What PC gets for free is the error
recursion. The gradient identity has to be bought, and the price is a schedule.

\section{Exact gradients live on a schedule --- and only on some graphs}
\label{sec:schedule}

\subsection{The schedule}

Song et al.'s Z-IL result \cite{song2020} recovers the exact gradient by timing rather
than by settling. Index layers by distance from the output, so that forward consistency
reads $w_i a_{i+1} = a_i$, let $\rho$ be the output residual, and let backprop's errors
be $\delta_0 = \rho$, $\delta_{i+1} = w_i \delta_i$, so that backprop's gradient for
$w_i$ is $\delta_i a_{i+1}$.

\begin{theorem}[Exactness on schedule]\label{thm:zil}
If layer $i$'s weights are read at inference step $t = i$ --- the step at which the error
front first reaches layer $i$, while its source latent still equals the forward
activation $a_{i+1}$ --- then the scheduled PC update for $w_i$ equals $\delta_i a_{i+1}$
exactly, for every $i < d$.
\end{theorem}

The hypothesis is doing all the work. Read the same parameter at free equilibrium
instead and you are back at Theorem~\ref{thm:gradne}. The identity is a property of
\emph{when} you look, not of the fixed point.

\subsection{Which graphs admit one}

Since exactness depends on reading each parameter exactly when its error arrives, the
question becomes: on which graphs does ``the first arrival'' carry the \emph{whole}
error? Give each node a rank, let $R$ be the head rank, and let the contribution
travelling along edge $e$ arrive at time $R - \mathrm{rank}(\mathrm{tgt}(e))$. Node $v$'s
first arrival is $\tau(v) = \min_{e\,:\,\mathrm{src}(e)=v}\bigl(R -
\mathrm{rank}(\mathrm{tgt}(e))\bigr)$. Call the schedule \emph{exact at $v$} when the
contributions that have arrived by $\tau(v)$ already sum to the full aggregate.

It pays to separate two things the literature runs together: whether a schedule
\emph{numerically} gets the right number, and whether it \emph{structurally} waited for
everything. Call a node \emph{admissible} at time $t$ when every contribution it needs
has arrived by $t$, and let the \emph{missed aggregate} be the signed sum of the
contributions that have not.

\begin{theorem}[Exactness, classified]\label{thm:signed}
For any ranked DAG, any arrival map and any per-node read time, the schedule is exact at
every node if and only if every node's missed aggregate is zero --- equivalently, iff
every node is either admissible or has its missed contributions exactly cancel. A
schedule can therefore be exact while being inadmissible, and that happens precisely
when some node's missed contributions cancel.
\end{theorem}

\leansignaturelink{lean/mettapedia/Mettapedia/MachineLearning/NeuralNetworks/%
PredictiveCoding/DAGScheduleExactness.lean}{218}
\begin{leansignature}
theorem dagScheduleExactAtEveryNode_iff_each_admissible_or_cancellation
    (G : SharedLatentDAG Node Edge) (parentError : Node → ℝ)
    (arrival : Edge → ℕ) (nodeTime : Node → ℕ) :
    dagScheduleExactAtEveryNode G parentError arrival nodeTime ↔
      ∀ v,
        dagScheduleAdmissible G arrival (nodeTime v) v ∨
          dagMissedContributionsCancel G parentError arrival (nodeTime v) v
\end{leansignature}

Cancellation is not hypothetical: a node missing contributions $+1$ and $-1$ is
inadmissible and numerically exact at once. So ``exact'' and ``waited for everything''
are genuinely different properties, and only the second is structural.

They fuse exactly when cancellation is impossible --- which is the regime the PC
literature actually works in, since parent contributions there are strictly positive.
Under that hypothesis exactness reduces to admissibility, and admissibility is decidable
by rank alone. The natural conjecture at that point is unit levelizability: every edge
steps one rank, $\mathrm{rank}(\mathrm{tgt}(e)) = \mathrm{rank}(\mathrm{src}(e)) + 1$, so
everything arrives in lockstep. That is sufficient. It is not necessary, and the true
criterion is weaker and more useful.

\begin{theorem}[Exact schedules, characterized]\label{thm:sched}
Let $G$ be a finite ranked DAG in which every parent contribution is strictly positive.
Then the first-arrival schedule is exact at every node if and only if
\[
\forall v,\ \forall e_1, e_2 \text{ with } \mathrm{src}(e_1)=\mathrm{src}(e_2)=v:\quad
\mathrm{rank}(\mathrm{tgt}(e_1)) = \mathrm{rank}(\mathrm{tgt}(e_2)).
\]
That is: exactly when no node's outgoing edges mix target ranks.
\end{theorem}

\leansignaturelink{lean/mettapedia/Mettapedia/MachineLearning/NeuralNetworks/%
PredictiveCoding/DAGScheduleExactness.lean}{727}
\begin{leansignature}
theorem dagFirstArrivalScheduleExact_iff_uniformOutgoingTargetRanks
    (G : SharedLatentDAG Node Edge) (parentError : Node → ℝ)
    (headRank : ℕ) (hhead : ∀ n, G.rank n ≤ headRank)
    (hpositive : ∀ e, 0 < dagParentContribution G parentError e) :
    dagFirstArrivalScheduleExact G parentError headRank ↔
      HasUniformOutgoingTargetRanks G
\end{leansignature}

Unit levelizability implies the criterion trivially --- all targets sit at
$\mathrm{rank}(v)+1$ --- but the converse fails, and we exhibit the witness: a graph with
uniform long skips is \emph{not} unit-levelized and \emph{is} exactly schedulable. Skips
are not the problem; \emph{ragged} skips are. A residual connection in which a parameter
feeds both a one-edge shortcut and a three-edge main path violates the criterion and is
proven inexact, which is the mechanism behind Salvatori et al.'s observation that the
construction is brittle on non-chain topologies \cite{salvatori2022}.

There is also a value-independent signed classification. Group every missed edge by its
target $u$, and let $g_{v,u}^{\rm miss}$ be the sum of its gains.

\begin{theorem}[Structural signed exactness]\label{thm:structsigned}
A schedule is exact for \emph{every} parent-error assignment if and only if
$g_{v,u}^{\rm miss}=0$ for every source $v$ and target $u$. The same theorem holds for
messages in any faithful real module. Thus signed skips may cancel structurally, but
only within a target group; with target-injective nonzero gains, or common-sign gains,
universal exactness reduces to admissibility.
\end{theorem}

The cancellation witness is therefore not a lucky choice of error values: two missed
edges to the same target, with gains $+1$ and $-1$, are universally exact. A single
missed nonzero edge is universally inexact. The scalar and vector-valued fixtures prove
both sides.

Read constructively, Theorem~\ref{thm:sched} is an architecture design rule:

\begin{quote}
Skip connections are compatible with exact backprop-equivalent PC scheduling precisely
when they do not mix target ranks. Uniform long skips are fine; rank-mixing residuals
are not.
\end{quote}

This is the one place where our theory tells an architect what to build rather than what
to avoid, and it says the space of exactly-schedulable graphs is strictly larger than
the field assumed.

\section{Error coordinates: what ePC preserves}
\label{sec:epc}

Goemaere et al.'s ePC \cite{goemaere2025} replaces latent states by prediction errors.
For arbitrary nonlinear scalar layers $f_i$, fixed input $x$, and errors $e_i$, define
recursively
\[
z_0=x,\qquad z_{i+1}=f_i(z_i)+e_i.
\]
Conversely, any input-clamped state stream gives
$e_i=z_{i+1}-f_i(z_i)$.

\begin{theorem}[Exact state/error reparameterization]\label{thm:epc}
These constructions are mutually inverse. On every finite prefix the ePC energy equals
the sPC energy after state reconstruction, and the corresponding local parameter update
is identical. At the differential level,
\[
d(E_{\rm sPC}\!\circ z)_e=(dE_{\rm sPC})_{z(e)}\circ dz_e.
\]
If $dz_e$ has a continuous-linear right inverse, ePC and sPC have exactly the same
critical-point condition.
\end{theorem}

This is a coordinate equivalence, not a blanket claim that every finite-step algorithm
has the same trajectory. The initialization rate still controls how far a signal has
travelled and how large it is when it arrives.

\begin{theorem}[First-arrival attenuation]\label{thm:wave}
For the one-edge-per-step recurrence $s_0=g$ and $s_{k+1}=\lambda s_k$,
\[
s_{L-i}=\lambda^{L-i}g.
\]
If $|\lambda|<1$, then for every positive tolerance the exact real signal eventually
falls below it. It need not equal zero: a machine-checked fixture is nonzero while below
tolerance. At $\lambda=1$, a unit signal never falls below any tolerance at most one.
A one-step ePC local update is $\lambda$ times the backprop product, and is exactly that
product at unit rate.
\end{theorem}

The tolerance statement is deliberately not called floating-point underflow. It is an
exact real-number decay theorem plus an explicit external detection policy.

\section{The cost of depth}
\label{sec:depth}

Settling is relaxation, and relaxation on a long clamped chain is slow. This is folklore;
here is the exact rate. On a path of $n$ segments with both ends clamped, Jacobi settling
is $(Jv)_k = \tfrac12(v_{k-1}+v_{k+1})$ at interior nodes and $0$ at the ends.

\begin{theorem}[Exact relaxation rate]\label{thm:depth}
The mode $s_k = \sin(k\pi/n)$ satisfies $Js = \cos(\pi/n)\, s$, hence
\[
(J^t s)_k \;=\; \cos(\pi/n)^t \, \sin(k\pi/n)
\qquad\text{for all } t .
\]
The spectral gap obeys $0 \le 1 - \cos(\pi/n) \le \pi^2/(2n^2)$, and
$\cos(\pi/n) \to 1$ as $n \to \infty$.
\end{theorem}

\leansignaturelink{lean/mettapedia/Mettapedia/MachineLearning/NeuralNetworks/%
PredictiveCoding/DepthPathology.lean}{106}
\begin{leansignature}
theorem pathSlowMode_iterate_exact
    (segments : ℕ) (hsegments : 0 < segments) :
    ∀ t n, n ≤ segments →
      pathJacobiIterate segments (pathSlowMode segments) t n =
        pathRelaxationRate segments ^ t * pathSlowMode segments n
\end{leansignature}

So the slowest mode decays by exactly $\cos(\pi/n)$ per sweep and the gap closes
quadratically in depth: doubling depth quarters the progress per settling step. This is
the mathematical spine under the engineering literature devoted to escaping it
\cite{goemaere2025}.

Real networks are not uniform, so the natural worry is whether one deep block poisons an
otherwise shallow architecture. It does.

\begin{theorem}[Heterogeneous blocks]\label{thm:het}
For a finite family of blocks, each with its own depth $n_i$, coupling eigenvalue $c_i$
and resulting exact mode rate, the family's modal spectral radius $\rho$ satisfies
$\rho < 1$ if and only if every block's rate is contractive; and if any block $i$ has
unit coupling, then
\[
1 - \rho \;\le\; \frac{\pi^2}{2\,n_i^{\,2}} .
\]
\end{theorem}

The bound is the pessimistic one: the family's spectral gap is capped by the
\emph{deepest} unit-coupled block, no matter how shallow everything else is. Stability
is a conjunction over blocks; speed is a minimum over blocks.

The corresponding escape is real, but only at an exact operator boundary.

\begin{theorem}[Frozen-backbone restriction boundary]\label{cor:adapter}
Join a backbone path of depth $B>0$ to an adapter path of depth $A>1$.  The resulting
length-$(B+A)$ block Jacobi operator is provably not block diagonal at the attachment.
Now freeze by embedding an adapter state into that full operator, applying the full
step, and projecting while deleting the clamped attachment row.  The restricted
operator is exactly the depth-$A$ adapter Jacobi operator, and a scalar $r$ is its
first-sine-mode multiplier if and only if
\[
r=\cos(\pi/A).
\]
Consequently its first-mode gap is independent of $B$ and obeys
$1-r\le\pi^2/(2A^2)$.  If the backbone instead remains active as a unit-coupled block,
then the global modal gap again obeys the backbone bound
$1-\rho\le\pi^2/(2B^2)$.
\end{theorem}

\leansignaturelink{lean/mettapedia/Mettapedia/MachineLearning/NeuralNetworks/%
PredictiveCoding/FrozenAdapterRestriction.lean}{233}
\begin{leansignature}
theorem isFrozenAdapterSlowModeRate_iff
    (backboneSegments adapterSegments : ℕ)
    (hadapter : 1 < adapterSegments) (rate : ℝ) :
    IsFrozenAdapterSlowModeRate backboneSegments adapterSegments rate ↔
      rate = pathRelaxationRate adapterSegments
\end{leansignature}

The independence is derived from restricting the coupled operator, rather than stored
in an adapter definition.  It is also sharp about the idealization: a depth-three
backbone with a depth-two adapter has frozen adapter rate $0$, but modal radius $1/2$
when both unit-coupled paths remain active.  Thus shallow adapters escape only when
freezing really deletes the backbone coordinates from settling; partial freezing or a
shared active relaxation can re-import the backbone depth.

The same first-arrival theorem also prices the imbalance that Depth-$\mu$P was designed
to mitigate \cite{innocentimupc2025}.  Let layer $i$ be $L-i$ links from the output and
let every inference link transmit at rate $\lambda$.  Composing the sealed signal,
one-step-update, and loss theorems gives the following result.

\begin{theorem}[Depth-credit boundary]\label{thm:mupcdepth}
The effective per-layer preconditioner is
\[
P_i=\lambda^{L-i},
\qquad
\ell_i^+=(1-\lambda^{L-i})^2\ell_i.
\]
For heterogeneous rates it is the suffix product $P_i=\prod_{j>i}\lambda_j$.  Equalizing
all $P_i$ forces every used rate to equal one, because the output's empty product is one.
Moreover, for $0<\lambda<1$, multiplying by any fixed polynomial in depth still tends to
zero; exact cancellation requires the exponential reciprocal $\lambda^{-(L-i)}$.
\end{theorem}

\leansignaturelink{lean/mettapedia/Mettapedia/MachineLearning/NeuralNetworks/%
PredictiveCoding/MuPCDepthCredit.lean}{366}
\begin{leansignature}
theorem muPCDepthScaling_cannot_cancel_contractingDepthCredit
    (width depth : ℕ) (inferenceRate : ℝ)
    (hwidth : 0 < width) (hdepth : 0 < depth)
    (hrateNonzero : inferenceRate ≠ 0) (hrateContractive : |inferenceRate| < 1) :
    depthMuPAdamLearningRateScale width depth *
        depthMuPHiddenMultiplier width depth = 1 ∧
      (∀ distance,
        depthMuPAdamLearningRateScale width depth *
            depthMuPHiddenMultiplier width depth *
              inferenceRate ^ distance = inferenceRate ^ distance) ∧
      (∀ learningRate distance,
        learningRate * inferenceRate ^ distance = 1 ↔
          learningRate = (inferenceRate ^ distance)⁻¹) ∧
      (∀ coefficient degree,
        ∀ᶠ distance : ℕ in atTop,
          polynomialDepthCompensation coefficient degree distance *
              inferenceRate ^ distance ≠ 1)
\end{leansignature}

The Depth-$\mu$P multiplier and its prescribed Adam learning-rate factor cancel each
other algebraically, but neither changes the inference rate.  Thus they repair a width
parameterization, not this exponential credit factor; indeed the reported $\mu$PC
experiments omit that Adam rescaling because it was empirically untrainable
\cite{innocentimupc2025}.  This is a scoped result, not a claim that $\mu$PC has no value:
it identifies exactly which instability its parameterization cannot remove.  The same
module also resolves a nearby verbal dispute about whether equilibrated activities
``barely move'': in a uniform-precision unit chain, node $k$ moves by exactly
$(k/L)(y-x)$.  Early fixed nodes move little as depth grows, while nodes near the output
need not.

\section{What a preconditioner can buy}
\label{sec:precond}

The standard response to slow or misbehaving settling is to rescale the update ---
precision scaling, natural-gradient-style metrics, or a line search. In the exactly
solvable case, what rescaling buys is a single number, and the answer is unexpectedly
sharp.

Take one chain link: source activation $s$, downstream gain $g$, weight $w$, prediction
$g w s$, target $y$, residual $r(w) = gws - y$ and loss $\ell(w) = \tfrac12 r(w)^2$.
Write $\gamma = gs$ for the effective gain and assume $\gamma \neq 0$. A PC update
preconditioned by $P$ moves along the normalized backprop direction,
$w^+ = w - P\,\ell'(w)/\gamma^2$.

\begin{theorem}[The $(P-1)^2$ law]\label{thm:precond}
$r(w^+) = (1-P)\,r(w)$, and therefore
\[
\ell(w^+) \;=\; (P-1)^2\, \ell(w).
\]
Consequently, for $r(w) \neq 0$:
\begin{enumerate}
\item $\ell(w^+) = 0$ \emph{if and only if} $P = 1$;
\item $\ell(w^+) < \ell(w)$ \emph{if and only if} $0 < P < 2$;
\item the shortfall against the ideal step is exactly
  $\Delta\ell(1) - \Delta\ell(P) = (P-1)^2\,\ell(w)$;
\item descent degrades strictly monotonically in the squared departure $(P-1)^2$.
\end{enumerate}
\end{theorem}

Every arm of our ignition ladder is a point on this parabola. An update in the correct
direction at $0.5\%$ magnitude is $P \approx 0.005$: it descends --- it is inside
$(0,2)$ --- while retaining $(0.995)^2 \approx 99\%$ of the loss per step. That is not a
tuning failure, it is the law. And overcorrecting is worse than useless: $P = 3$ lies
outside the descent window and \emph{quadruples} the loss, which is what our
precision-scaled arm did before it diverged. The window is $0 < P < 2$ and the target is
$P = 1$, with quadratic penalty for missing.

The vector case is where the one-step picture must not be over-read. If the residual is
an eigenvector of the preconditioner with eigenvalue $\mu$, the one-step law applies with
$\mu$ in place of $P$: $\ell \mapsto (\mu-1)^2 \ell$. It is tempting to conclude that a
preconditioner is correct only where its spectrum sits at $1$. That reading is wrong, and
the asymptotic truth is more permissive.

\begin{theorem}[Convergence is a spectral-radius condition]\label{thm:spectral}
Write $D = I - P$ for the preconditioner's departure from the identity. In any complete
normed algebra, $\rho(D) < 1$ if and only if $D^n \to 0$; hence the preconditioned
residual iterate converges to zero exactly when $\rho(I-P) < 1$. No diagonalizability is
assumed, so nonnormal operators are covered.
\end{theorem}

\leansignaturelink{lean/mettapedia/Mettapedia/MachineLearning/NeuralNetworks/%
PredictiveCoding/OperatorStability.lean}{22}
\begin{leansignature}
theorem spectralRadius_lt_one_iff_pow_tendsto_zero
    {A : Type*} [NormedRing A] [NormedAlgebra ℂ A] [CompleteSpace A]
    [NormOneClass A] [Nontrivial A] (a : A) :
    spectralRadius ℂ a < 1 ↔
      Tendsto (fun n : ℕ => a ^ n) atTop (𝓝 0)
\end{leansignature}

So $P = I$ is the \emph{one-step optimum}, not the requirement. Asymptotically an entire
open region works, and one may sit far from the identity on every mode and still
converge. For a scalar this is the familiar $|1-P| < 1$ --- the window $0 < P < 2$ of
Theorem~\ref{thm:precond} --- and the spectral form is what survives to matrices.

Finite-time behavior needs its own theorem.

\begin{theorem}[Geometric envelope]\label{thm:envelope}
If $\rho(D)<q<1$, then there is a finite $C\geq1$ such that
\[
\|D^n\|\leq Cq^n\quad\text{and}\quad
\|D^nr\|\leq Cq^n\|r\|\qquad(n\geq0).
\]
The constant absorbs the finite nonnormal transient; it cannot in general be determined
from the spectral radius alone.
\end{theorem}

The obstruction is exact for the Jordan family
\[
D_{\lambda,c}=\begin{pmatrix}\lambda&c\\0&\lambda\end{pmatrix},\qquad
D_{\lambda,c}^{,n}=
\begin{pmatrix}
\lambda^n&nc\lambda^{n-1}\\0&\lambda^n
\end{pmatrix}.
\]

\begin{theorem}[Arbitrarily large stable transients]\label{thm:jordan}
$D_{\lambda,c}$ has spectrum $\{\lambda\}$, and is nonnormal whenever $c\ne0$.
For $\lambda=0$ it satisfies $D^2=0$, yet a unit residual is amplified on the first
step by exactly $|c|$. Consequently, for every proposed bound $B\geq0$ there is a
nonnormal operator with spectral radius zero, exact two-step convergence, and first-step
amplification greater than $B$.
\end{theorem}

The earlier witness $c=2$ is therefore one point in an unbounded family, not a curiosity.
Asymptotic stability, geometric rate, transient size, and step-wise descent are four
different claims. A nonnormal preconditioner may honour the first two while violating
the fourth, with the third controlled only by additional operator information.

What no amount of spectral latitude buys is a single global knob.

\begin{theorem}[One preconditioner cannot serve two modes]\label{thm:mode}
There is a two-mode fixture and a single diagonal preconditioner
$\mathrm{diag}(\tfrac12, 3)$ that strictly \emph{decreases} the loss on one eigenmode and
strictly \emph{increases} it on the other, from equal starting losses.
\end{theorem}

This is the no-go for scaling one's way out: a metric that is merely large, or merely
uniform, buys nothing, because its eigenvalues are wrong in both directions at once. The
bar is spectral rather than magnitudinal --- but by Theorem~\ref{thm:spectral} it is a bar
of membership, $\rho(I-P)<1$, not of identity. Information-geometric preconditioning
therefore keeps an honest target that is not merely ``be a second-order method''.

Line search fares better, because it targets $P$ adaptively rather than fixing it.

\begin{theorem}[Backtracking descends]\label{thm:armijo}
Under an explicit smoothness bound, Armijo backtracking terminates --- some
$P = P_0\sigma^n$ satisfies the Armijo condition --- and every accepted step is
non-increasing in energy. This instantiates with the exact smoothness constant of the
depth-2 PC energy.
\end{theorem}

This is the correctness certificate for the settling procedure used by the surviving
experimental arms, whose logged energies were monotone in every accepted run.

\section{Not every equilibrated-energy saddle is strict}
\label{sec:saddles}

Innocenti et al.\ prove strictness for important classes of deep-linear equilibrated PC
energy and conjecture that every remaining saddle is strict \cite{innocenti2025}. A
strict saddle is a critical point with a direction of negative second curvature. The
unrestricted conjecture has a covariance-degenerate boundary that the partial theorem
does not cover.

Take a scalar two-layer model and the two examples with inputs $(1,1)$ and targets
$(1,-1)$. Its exact equilibrated energy at weights $(a,b)$ is
\[
F(a,b)=\frac{1+a^2b^2}{2(1+b^2)}.
\]
This formula is derived in Lean from the rescaled batch residual objective, rather than
postulated. Every $(a,0)$ is critical because the input--target covariance vanishes.

\begin{theorem}[Strict region and non-strict boundary]\label{thm:saddles}
At $(a,0)$, the output-weight direction has second curvature $a^2-1$. Hence every
$a^2<1$ has a strict-saddle certificate, including the published origin case. But at
$(1,0)$ every directional second curvature is zero, while for every $0<t<1$,
\[
F(1-t,t)<F(1,0)<F(1+t,t).
\]
Thus $(1,0)$ is a genuine saddle with explicit descending and ascending approach curves,
but it is not strict.
\end{theorem}

\leansignaturelink{lean/mettapedia/Mettapedia/MachineLearning/NeuralNetworks/%
PredictiveCoding/EquilibratedEnergySaddles.lean}{319}
\begin{leansignature}
theorem orthogonalDataset_universalStrictSaddleConjecture_counterexample :
    IsCriticalPoint orthogonalDatasetEquilibratedEnergy (1, 0) ∧
      HasTwoSidedSaddleCurves orthogonalDatasetEquilibratedEnergy (1, 0)
        (fun step => (1 - step, step))
        (fun step => (1 + step, step)) ∧
      (∀ direction,
        HasSecondDirectionalCurvatureAt orthogonalDatasetEquilibratedEnergy
          (1, 0) direction 0) ∧
      ¬ HasStrictSaddleCertificate orthogonalDatasetEquilibratedEnergy (1, 0)
\end{leansignature}

So the universal conjecture is false without a nondegeneracy hypothesis. This does not
contradict the published strictification theorem: the counterexample lies outside its
covered class. It identifies the missing issue more sharply --- a cubic-order saddle can
survive exactly where the quadratic curvature vanishes.

\section{The coupled system}

Settling and learning run together, and their interaction has its own stability region.
Write the two-timescale linear system with settling precision $\lambda$ and learning rate
$\eta$: the settled error is multiplied by $1-\lambda$ per step, and the weight by
$1-\eta\lambda$.

\begin{theorem}[Stability region]\label{thm:coupled}
Both multipliers are contractive if and only if
\[
0 < \lambda < 2 \qquad\text{and}\qquad 0 < \eta\lambda < 2 ,
\]
in which case the coupled iteration converges to zero from every initial state.
If $\eta\lambda > 2$, the learning eigenmode diverges in norm.
\end{theorem}

Note the shape: $0 < \eta\lambda < 2$ is the same window as $0 < P < 2$ in
Theorem~\ref{thm:precond}, which is not a coincidence --- both are the condition for a
one-step affine contraction. It is also the formal account of the late-run divergence one
of our arms exhibited: the product $\eta\lambda$, not either factor, is what must stay
inside the window.

The second dynamical fact concerns what settling does to the \emph{shape} of the output
distribution, and it is the one our experiments measured directly.

\begin{theorem}[Entropy turnover]\label{thm:turnover}
For a linear signal of any depth $d > 0$ and any soft target $p \in (0,1)$ with
$p \neq \tfrac12$, the settled output's entropy $H(s)$ after $s$ settling steps satisfies
\[
H(0) > H_b(p), \qquad H(d) = H_b(p), \qquad H(d+1) > H_b(p),
\]
where $H_b$ is binary entropy. Hence $H$ is neither monotone nor antitone; and it is
non-monotone \emph{if and only if} the target is non-uniform. A gradient-trained
comparator provably lands strictly below the target entropy.
\end{theorem}

Under-settle and the output is diffuse; settle exactly and you hit the target; over-settle
and it re-diffuses past it. Settling depth is therefore not a ``more is better'' knob:
there is an interior optimum, and it is at the signal depth. Our tree-network campaign
observed this curve --- including the turnover --- before we proved it.

The turnover is not an artifact of linearity. Let the settling map be any monotone
contraction fixing zero, with rate $q < 1$ --- a class that includes the saturating
nonlinearities PC implementations actually use.

\begin{theorem}[Nonlinear turnover, sufficient conditions]\label{thm:nonlin}
Fix such a settling map, a depth $d>0$, and a non-uniform target $p \in (0,1)$,
$p \neq \tfrac12$. If the depth-$d$ iterate reaches the target logit from the seed, then
the same three regimes hold --- entropy above target before arrival, equal at depth,
above again afterwards --- and the entropy is neither monotone nor antitone. Moreover
contraction bounds the reachable set: $|\operatorname{logit}(p)| \le q^{\,d}\,|\text{seed}|$.
\end{theorem}

The reachability hypothesis is doing real work and we flag it rather than bury it: this is
a theorem about a class of nonlinear settling dynamics that \emph{can} hit their target,
not a theorem about arbitrary trained nonlinear PC networks. The contraction bound is the
honest sting --- it says which targets a given depth can reach at all. We therefore build
and train one concrete nonlinear model rather than assume reachability.

Let $\phi(x)=\sigma(x)-\tfrac12$. This centered sigmoid fixes zero, is strictly
increasing and globally $\tfrac14$-Lipschitz, and is proved not to agree pointwise with
any real-linear map. Put a parameter $w$ one edge upstream, read
$p_w=\sigma(\phi(w))$ after one settling step, choose $p_1$ as the training target, and
train with $L(w)=(p_w-p_1)^2$.

\begin{theorem}[Trained nonlinear turnover]\label{thm:trainednonlin}
$w=1$ is the unique zero-loss parameter and a global minimizer. At that trained weight,
the output probabilities at settling depths $0,1,2$ are
$\tfrac12,p_1,\tfrac12$, with $p_1>\tfrac12$, so entropy strictly decreases and then
increases. The middle entropy remains strictly below uniform throughout an explicit
open probability neighborhood of $p_1$. The dynamics lift coordinatewise to every
finite vector width; an all-one seed gives turnover in every coordinate, while an
all-zero seed stays uniform.
\end{theorem}

\leansignaturelink{lean/mettapedia/Mettapedia/MachineLearning/NeuralNetworks/%
PredictiveCoding/TrainedNonlinearTurnover.lean}{432}
\begin{leansignature}
theorem trainedNonlinearEntropyTurnover_crown :
    (¬ ∃ L : ℝ →ₗ[ℝ] ℝ, ∀ x, L x = centeredSigmoidActivation x) ∧
      trainedCenteredSigmoidProblem.loss 1 = 0 ∧
      (∀ weight, trainedCenteredSigmoidProblem.loss 1 ≤
        trainedCenteredSigmoidProblem.loss weight) ∧
      Real.binEntropy trainedCenteredSigmoidTargetProbability <
        trainedCenteredSigmoidProblem.entropy 1 0 ∧
      trainedCenteredSigmoidProblem.entropy 1 1 =
        Real.binEntropy trainedCenteredSigmoidTargetProbability ∧
      Real.binEntropy trainedCenteredSigmoidTargetProbability <
        trainedCenteredSigmoidProblem.entropy 1 2 ∧
      (¬ ∃ seed : ℝ, centeredSigmoidSettling.activation seed = 1) ∧
      0 < trainedCenteredSigmoidRobustRadius
\end{leansignature}

This separates three notions that ``contractive nonlinear PC'' can obscure. Contraction
controls sensitivity, training certifies the chosen target, and the one-edge delay causes
turnover. None implies the others: the same centered sigmoid is contractive but its range
is $(-\tfrac12,\tfrac12)$, so no seed can reach logit $1$.

\section{Fit does not imply guidance}
\label{sec:yield}

The last result is not about PC at all. It is about what PC-versus-BP comparisons
measure, and it is the reason our empirical verdict is stated in the currency it is.

A search-guidance prior $\theta$ over $k$ program tokens is sampled i.i.d.\ within a
budget $B$; a \emph{hit} is a draw of the target.

\begin{theorem}[Yield is target mass]\label{thm:yield}
The expected number of target hits is exactly $B\,\theta_{\mathrm{target}}$, and is
strictly increasing in $\theta_{\mathrm{target}}$.
\end{theorem}

That is unsurprising, and it is the setup for the separation.

\begin{theorem}[Equal loss, ordered yield]\label{thm:sep}
Let $\theta^{\mathrm{lo}} = (\tfrac34,\tfrac14)$ and
$\theta^{\mathrm{hi}} = (\tfrac14,\tfrac34)$, with outcome $1$ the target. Against the
soft target $(\tfrac12,\tfrac12)$ these priors have \emph{equal} cross-entropy, yet for
every budget $B > 0$ their expected yields are strictly ordered:
\[
\tfrac{B}{4} \;<\; \tfrac{3B}{4}.
\]
\end{theorem}

\leansignaturelink{lean/mettapedia/Mettapedia/MachineLearning/SearchGuidance/%
DecisivenessYield.lean}{242}
\begin{leansignature}
theorem equal_crossEntropy_strictly_ordered_expectedYield :
    softHalfTargetCrossEntropy lowTargetMassPrior =
        softHalfTargetCrossEntropy highTargetMassPrior ∧
      ∀ budget : ℕ, 0 < budget →
        iidExpectedTargetHits lowTargetMassPrior budget 1 <
          iidExpectedTargetHits highTargetMassPrior budget 1
\end{leansignature}

Two models, identical loss, threefold difference in search performance. Cross-entropy is
symmetric under swapping mass between outcomes; yield is not, because only one outcome is
the target. So training loss --- the quantity every PC-versus-BP comparison in the
literature reports, including ours --- provably does not determine search performance.

What does determine it, in this model, is the mass the prior puts on the target: its
\emph{decisiveness}. This is why Section~\ref{sec:tree} reports prior entropy alongside
loss, and why the entropies are the numbers that moved with performance.

\section{The experiments}
\label{sec:experiments}

\subsection{Paired self-learning campaign on the tree network}
\label{sec:tree}

We reproduced Gauthier's OEIS program-synthesis pipeline byte-faithfully (Standard ML,
HOL4, tree neural network in C) and implanted PC \emph{in place} in the C training code,
guarded by numerical oracles (analytic golden vectors to $10^{-9}$--$10^{-16}$; the BP
path byte-identical when PC is off). The learning rule is the only variable.

Single-generation results establish the mechanism. PC's yield climbs with settling depth
--- $63\%, 67\%, 83\%, 93\%$ of BP's newly solved sequences at $4, 8, 16, 32$ settling
steps --- at roughly linear cost, up to $\sim$10$\times$ BP's training time at 32 steps,
and then \emph{turns over} at 64 steps. Throughout, the entropy of the search-guidance
prior tracked yield at every point, including the turnover, while training loss did not:
PC fit the training objective as well as or better than BP while guiding search worse.
Measured prior entropies: BP $0.671$; PC $1.479 \to 0.975$ over steps $4 \to 32$,
re-diffusing to $1.154$ at step 64. Theorem~\ref{thm:turnover} is the shape of that
curve and Theorem~\ref{thm:sep} is why loss failed to see it.

The multi-generation campaign makes the causal comparison real: three paired seed-worlds,
one frozen random generation-0 each, both arms grown for five self-training generations
under identical budgets. Cumulative solved OEIS sequences:

\begin{center}
\begin{tabular}{lrrrrr}
\toprule
 & gen-1 & gen-2 & gen-3 & gen-4 & gen-5 \\
\midrule
BP, world 1 & 8{,}752 & 11{,}709 & 13{,}224 & 14{,}566 & 15{,}852 \\
BP, world 2 & 8{,}361 & 11{,}291 & 13{,}250 & 15{,}054 & 16{,}081 \\
BP, world 3 & 8{,}627 & 11{,}510 & 13{,}267 & 14{,}793 & 15{,}843 \\
PC, world 1 & 6{,}465 & 6{,}914 & 9{,}067 & 9{,}486 & 10{,}211 \\
PC, world 2 & 6{,}093 & 8{,}031 & 8{,}359 & 9{,}195 & 9{,}411 \\
PC, world 3 & 7{,}548 & 7{,}548 & 8{,}820 & 10{,}111 & 10{,}296 \\
\midrule
PC/BP ratio & 78\% & 65\% & 66\% & 65\% & 62.6\% \\
\bottomrule
\end{tabular}
\end{center}

BP compounds metronomically: the three endpoints land within $1.5\%$ of one another
($15{,}843$--$16{,}081$). PC ignites but advances in erratic bursts
--- including one genuine zero-progress generation (world~3, generation~2, byte-identical
solution set) --- and finishes between $58.5\%$ and $65.0\%$ of BP's frontier (world~1
$64.4\%$, world~2 $58.5\%$, world~3 $65.0\%$), $62.6\%$ in aggregate. Same frozen start,
same search budget: the learning rule alone costs a third of the frontier, at
$32$ settling steps' training overhead.

\paragraph{Unique coverage: does PC find solutions BP does not?}
Yes, and the answer has a shape. Per world, PC solves $242/277/286$ sequences its paired
BP arm misses ($\approx 2.5\%$ of PC's own catalogue, against $\sim 6{,}000$ the other
way). Pooling all three worlds, PC's union contributes $246$ sequences that \emph{no} BP
world solved --- a combined catalogue of $19{,}280$ against BP's $19{,}034$. One case is
reproducible rather than stochastic: a single sequence (\texttt{A087299}) solved by PC in
all three independently evolved worlds, by BP in none, with the \emph{identical}
$11$-token program each time --- an existence proof that the two rules steer search into
genuinely different territory. The mechanism is visible in program lengths (first-listed
program per solved sequence, pooled over worlds): the sequences only BP solved take
median-$24$-token programs (mean $32.4$), the sequences only PC solved median $19$ (mean
$20.7$), and at the generation-five frontier the gap widens --- BP's newly solved
sequences run to median $40$ tokens (mean $47.3$) against PC's $25$ (mean $28.1$).
Decisive guidance sustains long derivations; flatter guidance spreads mass across short
ones and occasionally covers a short solution that confident-but-elsewhere mass walks
past. The economics, however, are
one-sided: adding a \emph{second BP world} contributed $2{,}241$ new sequences and a third
$941$, so as a diversity source PC is dominated several-fold by re-running BP on a fresh
seed --- on this topology. Whether richer topologies (residual adapters on a frozen
backbone --- the one regime the depth theorems endorse --- or graph-structured models)
change PC's standing is deliberately left open and would require a separately registered
comparison.

\paragraph{Campaign closure and provenance.}
All thirty phase artifacts (three worlds $\times$ two arms $\times$ five guided
generations) are sealed with per-file SHA-256 evidence manifests, and every number in
this subsection was recomputed from those sealed artifacts before publication. The
zero-progress generation is hash-verified: world~3's PC solution file carries the
identical digest at generations one and two. Four operational amendments were filed
during the run --- memory-floor recovery and driver-relaunch policy --- each sealed
against the campaign seal and each recording \emph{protocol unchanged, science
parameters unchanged}; the arms, the checker, and the validation path were never
touched. The campaign halted at generation five, the pre-declared mandatory guided
generation, by an operator decision recorded before the final arms completed; the
protocol's adaptive stopping rule (earliest adaptive stop at generation seven, hard stop
at ten) was consequently never evaluated, and no claim here rests on it. Extending the
same frozen worlds toward the hard stop remains available as a pre-registered fork.

\subsection{The seq2seq ignition ladder}
\label{sec:ladder}

On a faithful reimplementation of the alien-coding NMT baseline \cite{gauthier2023}
(14-token program language, 512-unit bidirectional-LSTM encoder--decoder with scaled-Luong
attention; the trained reference solves 1{,}860 new OEIS sequences on seed 1, with 222
target-associated hits against a randomized-association null of 6), we ran a seven-arm
overfit-32 ignition ladder for PC variants. The registered composite bar: final loss
$\le 0.05$; at least 30/32 exact greedy decodes; monotone settling energies; and
\emph{distinctness} --- the final updates must differ measurably from BP's, median
parameter cosine $< 0.95$ --- so that a pass cannot be secretly backpropagation. The BP
anchor passes at loss $0.026$, 32/32.

\begin{center}
\small
\begin{tabular}{llll}
\toprule
arm & result & loss / exact & measured failure mode \\
\midrule
fixed-rate ePC & FAIL & 105.4 / 0 & settling divergence \\
backtracking, 16-step & FAIL (time) & --- & exceeded registered wall-clock \\
backtracking, 5-step & FAIL & 11.2 / 1 & BP-direction updates at 0.5\% magnitude \\
precision-scaled & FAIL & 21.4 / 0 & late-run divergence \\
staged BP$\to$PC, 4000 & FAIL & 0.075 / 32 & loss bar only; distinctness passed \\
staged BP$\to$PC, 4500 & FAIL & 0.044 / 32 & distinctness only; updates BP-like \\
staged, deepened inference & \textbf{PASS} & 0.037 / 32 & all bars; median cosine 0.80 \\
\bottomrule
\end{tabular}
\end{center}

Three of these failures are theorems. The 5-step arm is $P \approx 0.005$ in
Theorem~\ref{thm:precond}: correct direction, $99\%$ of the loss retained per step. The
precision-scaled arm left the $0 < \eta\lambda < 2$ window of
Theorem~\ref{thm:coupled}. The fixed-rate arm is the same failure without the line search
that Theorem~\ref{thm:armijo} certifies.

The joint reading of the two staged arms is the sharpest empirical fact in this note: the
staged schedule that starts at provable BP-equivalence and walks toward PC \emph{learns
to the bar exactly as far as its updates remain BP-like}, and retains PC-distinctness
exactly when it misses the bar. The seventh arm broke the tension: continuing the staged
run with deepened inference passed every registered bar at once --- loss $0.0366$, 32/32
exact decodes, monotone settling, genuinely PC-distinct updates (median parameter cosine
$0.80$, 13/18 tensors below $0.95$).

That pass was seed~0 and judged at one terminal snapshot, so a protocol weakness had to be
settled before it could be believed: the terminal cosine statistic turns out to swing
\emph{within a single fixed stage} --- measured medians $0.929, 0.952, 0.322, 0.787,
0.801$ at updates $4501$ through $4700$ --- so no one endpoint is a robust mechanism
criterion. The replication was therefore pre-registered on an independent seed and on a
prospectively fixed five-checkpoint window instead, with the bars unchanged.

It passes. On seed~1, the complete staged schedule (3000 BP-equivalent one-step anchor
updates, then fixed transitions to five-step inference, 4700 updates in total) meets every
registered bar: final loss $0.0014$, 32/32 exact decodes, all 96 logged inference records
monotone, and mechanism distinctness at \emph{all five} checkpoints --- median of the
checkpoint medians $0.228$, against a required cutoff of $0.95$. The updates are thus far
\emph{more} PC-distinct than seed~0's terminal $0.80$, and distinct throughout the window
rather than at a chosen point. Staged ePC ignition reproduces.

Two scope limits are registered with it and remain: the schedule begins from a
BP-equivalent anchor, so this is not learning from random initialization; and it is
overfit-32 scale, not full-scale NMT performance.

\section{Verdict}
\label{sec:verdict}

\textbf{Dead, on the evidence:} predictive coding as a drop-in replacement for
backpropagation on synchronous digital hardware, for training program-synthesis models
from scratch. At small scale it is dominated at equal epochs and vastly dominated at
equal wall-clock; at medium multi-generation scale it holds a stable two-thirds frontier
deficit with erratic self-training dynamics; on the recurrent architecture, seven
registered arms produced exactly one passing design, and it is staged from a
BP-equivalent anchor rather than trained from random initialization.

\textbf{Surviving, with formal or measured support:}
\begin{enumerate}
\item \emph{Shallow PC adapters under an exact frozen restriction.} The depth pathology
  depends only on the adapter after the coupled operator is genuinely restricted
  (Theorem~\ref{cor:adapter}), and the empirical verdict independently points at this
  regime.  The theorem does not cover a shared active normalization or partial freezing:
  its negative half proves that an active unit-coupled backbone re-imports the backbone
  depth bound.
\item \emph{BP-train, PC-fine-tune.} The staged homotopy arms were the only ones to reach
  32/32 exact decodes; the deepened-inference continuation passed every bar on seed~0, and
  the pre-registered independent-seed replication passed on seed~1 with distinctness at
  every checkpoint of a fixed window. This is the one place PC ignited reproducibly under
  criteria fixed in advance. Buying PC's properties \emph{after} BP's learning is the
  empirically supported direction.
\item \emph{Second-order geometry, with the bar set by Theorem~\ref{thm:mode}.} The
  preconditioning failures are exactly characterized, so the cure is exactly specified:
  a metric whose departure from the identity is a contraction, $\rho(I-P)<1$, on the modes
  that matter --- a membership condition, not identity with backprop. Merely large, or
  merely uniform, is proven not to work.
\item \emph{Rank-uniform skip architectures.} Theorem~\ref{thm:sched} says exact PC
  scheduling survives skip connections that do not mix target ranks --- a strictly larger
  design space than levelizability, and untested.
\item \emph{Asynchronous and neuromorphic settings.} Untested here; nothing in this note
  bears on them, except that the decisiveness instrument of Section~\ref{sec:yield}
  travels to any setting where a learned prior guides search.
\end{enumerate}

\section{Limitations and open problems}

The models are exact and, where it matters, small: scalar paths, finite priors, expected
rather than at-least-one hit counts. Theorem~\ref{thm:sched} assumes strictly positive
parent contributions; Theorem~\ref{thm:signed} now covers the signed case, so what remains
open there is not the classification --- Theorem~\ref{thm:structsigned} supplies that ---
but an architecture-level criterion for useful cancellations under realistic learned
gains. Theorem~\ref{thm:nonlin} still needs reachability in its general form;
Theorem~\ref{thm:trainednonlin} discharges it for one trained centered-sigmoid model, not
for arbitrary coupled multivariate training. Theorem~\ref{thm:envelope} gives finite-time
nonnormal bounds but its constant is existential; sharp computable constants require more
operator structure. Theorem~\ref{thm:saddles} refutes universal strictness and classifies
one critical line, not every equilibrated-energy critical manifold. The ePC equivalence
is mathematical over exact reals; implementation-level rounding behavior remains a
separate question. The ignition ladder is overfit-32 scale: its purpose is
correctness-gated variant triage, not benchmark performance. The distinctness criterion's
placement at the registered final step is a design choice whose alternatives remain
untested. All ``dead'' verdicts are scoped to the regimes actually measured.

\appendix

\section{The causal boundary: clamping is not intervention}

\subsection{Observation is direction-blind}

Take centered variables with positive-definite covariance
\[
\Sigma=\begin{pmatrix}v_X&c\\c&v_Y\end{pmatrix},
\qquad \Delta=v_Xv_Y-c^2>0.
\]
The fitted $X\to Y$ quadratic PC model has coefficient $c/v_X$ and residual variance
$\Delta/v_X$; the reverse model has coefficient $c/v_Y$ and residual variance
$\Delta/v_Y$.  Their residual operators are provably distinct.  Nevertheless, direct
algebra reduces both full energies, pointwise for every real pair $(x,y)$, to
\[
\frac{v_Yx^2-2cxy+v_Xy^2}{\Delta}.
\]
The equality therefore lifts to every finite batch of \emph{arbitrary} observations.
There is no Gaussian data hypothesis: ``Gaussian'' names the quadratic energy.  More
generally, the formal batch theorem is
\[
\sum_s E(z_s)=\operatorname{tr}(Q\widehat\Sigma),
\qquad
Q=A^\top\Lambda A,
\quad
\widehat\Sigma=\sum_s z_s z_s^\top .
\]
Changing the data distribution while retaining this energy cannot expose higher-order
orientation information.  A non-quadratic residual log-density would change the model's
architecture, not merely its environment.

Ordinary PC clamping does not help.  Clamp $X=x$ but retain every residual term and
minimize over $Y$.  In both orientations the unique equilibrium is
$Y=(c/v_X)x$, the ordinary conditional mean.  PC clamping is conditioning, and
conditioning is direction-blind.

\subsection{One uniform severing operation}

Pearl's intervention replaces the intervened structural mechanism rather than
conditioning on its output \cite{pearl2012,halpern2000}.  At the PC energy level this is
one orientation-independent operation: clamp a node and delete the residual term owned
by that node.  Residual rows are indexed by their owning node, so the operation itself
never inspects edge direction.

\begin{theorem}[Intervention certificate]\label{thm:intervention}
Under $do(X=x)$, the $X\to Y$ equilibrium remains $(c/v_X)x$, whereas the $Y\to X$
equilibrium is $0$.  They differ if and only if $c\ne0$ and $x\ne0$.  The two
interventional energy functions are
\[
E^{do}_{X\to Y}(x,y)=\frac{(y-(c/v_X)x)^2}{2(\Delta/v_X)},
\qquad
E^{do}_{Y\to X}(x,y)=\frac{y^2}{2v_Y},
\]
and some intervention separates them if and only if $c\ne0$.
\end{theorem}

\leansignaturelink{lean/mettapedia/Mettapedia/MachineLearning/NeuralNetworks/%
PredictiveCoding/CausalDirectionIdentifiability.lean}{866}
\begin{leansignature}
theorem observationally_blind_intervention_certificate
    {Sample : Type*} [Fintype Sample]
    (covariance : PositiveBivariateCovariance)
    (hcorrelated : covariance.covXY ≠ 0) (x y : Sample → ℝ) :
    covariance.forwardResidualMatrix ≠ covariance.reverseResidualMatrix ∧
      (forwardNodeResidualModel covariance).residualMatrix ≠
        (reverseNodeResidualModel covariance).residualMatrix ∧
      (∑ sample, covariance.forwardGaussianPCEnergy (x sample) (y sample)) =
        ∑ sample, covariance.reverseGaussianPCEnergy (x sample) (y sample) ∧
      covariance.SeparatesDirectionsByInterventionalEnergy
\end{leansignature}

At the concrete covariance $\left(\begin{smallmatrix}1&1\\1&2\end{smallmatrix}\right)$
and $x=y=1$, both half-normalized observational energies equal $1/2$, while the
interventional energies are respectively $0$ and $1/4$.  The negative fixture is just
as important: independent equal-variance variables give $y^2/2$ in both orientations,
so intervention does not orient an absent edge.  This certificate supplies one bit about
one edge; it is not a general causal-discovery theorem.

An energy difference at a selected state is not yet a learning theorem, so we also let
both orientations fit.  On a finite intervention design $(x_s)_s$, take responses from
the true forward mechanism, $y_s=(c/v_X)x_s$.  A forward candidate fits its own
coefficient $\beta$ and a reverse candidate fits its own coefficient $\gamma$; both are
scored by the same severing operation.  Their batch objectives reduce to
\[
L_{X\to Y}(\beta)
 = \frac{(c/v_X-\beta)^2}{2(\Delta/v_X)}\sum_s x_s^2,
\qquad
L_{Y\to X}(\gamma)
 = \frac{(c/v_X)^2}{2v_Y}\sum_s x_s^2.
\]
The reverse candidates are structurally distinct as $\gamma$ varies, but their fitted
row is precisely the row deleted by $do(X)$; the second objective's independence from
$\gamma$ is therefore a theorem about the uniform intervention, not a restriction of
the competitor.

\begin{theorem}[Fitted two-node causal selection]\label{thm:causalselection}
For any finite design with $\sum_s x_s^2>0$, both orientation-specific objectives attain
their minima.  The minimized forward energy is strictly lower than the minimized reverse
energy if and only if $c\ne0$.
\end{theorem}

\leansignaturelink{lean/mettapedia/Mettapedia/MachineLearning/NeuralNetworks/%
PredictiveCoding/CausalDirectionSelection.lean}{316}
\begin{leansignature}
theorem forwardOrientationWinsByMinimizedInterventionalEnergy_iff
    {Sample : Type*} [Fintype Sample]
    (covariance : PositiveBivariateCovariance) (x : Sample → ℝ)
    (hexcited : 0 < ∑ sample, x sample ^ 2) :
    covariance.ForwardOrientationWinsByMinimizedInterventionalEnergy x ↔
      covariance.covXY ≠ 0
\end{leansignature}

Indeed, $\beta=c/v_X$ gives forward minimum zero, while every $\gamma$ minimizes the
reverse objective at the displayed positive value exactly when $c\ne0$.  At independent
equal variance both minima are zero, so selection fails.  The conclusion is strictly
two-node and assumes interventional responses from the forward linear mechanism,
positive residual scales, and an excited finite design; it is not a graph-discovery
claim.

\subsection{A separate continual-learning bridge}

Another causal-coding proposal turns on whether learning one thing and then another
is the same as learning both. It is not, and the defect is exactly computable in the
quadratic case. For two quadratic tasks with curvatures $A_1, A_2$, step size $\eta$, and
gradient $g_1$:
\[
\theta_{1\to2} - \theta_{\mathrm{add}} \;=\; \eta^2 A_2\, g_1(\theta),
\qquad
\theta_{1\to2} - \theta_{2\to1} \;=\; \eta^2 [A_1, A_2]\, \theta ,
\]
the second for centred tasks. So sequential learning departs from additive learning by the
second task's curvature acting on the first task's gradient, and the order in which you
learn matters exactly to the extent that the curvatures fail to commute.

These are two different phenomena, and we separate them rather than let one stand in for
the other. A fixture with oblique tasks has a non-zero commutator and is order-dependent.
A fixture in which one task is repeated has a \emph{zero} commutator --- no order
interference is possible --- and is still non-additive, because the same cause is counted
twice. On the evidence side that double counting is exact: re-applying a Gaussian
contribution yields the correct once-each ledger plus precisely one extra copy of that
contribution.

We state the scope plainly, and the formalization enforces it: what is proven is linear
and quadratic. The general correspondence between predictive coding and causal coding
remains explicitly conjectural beyond these instances, and its status now records that a
causal use requires an interventional clamp.  No theorem promotes the general
correspondence.

\section*{Where to find these}

All statements are machine-checked in Lean~4, in the Mettapedia library under the
namespace \leanname{Mettapedia.MachineLearning}. Each is axiom-clean, depending only on
\leanname{propext}, \leanname{Classical.choice} and \leanname{Quot.sound}. Files below sit
in \leanname{NeuralNetworks/PredictiveCoding/} unless marked otherwise.

{\small
\begin{longtable}{@{}l>{\raggedright\arraybackslash}p{0.61\textwidth}>{\raggedright\arraybackslash}p{0.205\textwidth}@{}}
\caption{Index of formal statements.}\label{tab:index}\\
\toprule
& Lean name & File \\
\midrule
\endfirsthead
\toprule
& Lean name & File \\
\midrule
\endhead
Thm~\ref{thm:bayes}
  & \leansource{lean/mettapedia/Mettapedia/MachineLearning/NeuralNetworks/PredictiveCoding/LinearGaussianOperator.lean}{219}{LinearGaussianOperatorModel.equilibrium\_iff\_eq\_posteriorMean}
  & LinearGaussian\-Operator \\
  & \leansource{lean/mettapedia/Mettapedia/MachineLearning/NeuralNetworks/PredictiveCoding/MatrixGaussianChain.lean}{452}{matrixPCEquilibrium\_iff\_eq\_conditionalPosteriorMean} & MatrixGaussian\-Chain \\
  & \leansource{lean/mettapedia/Mettapedia/MachineLearning/NeuralNetworks/PredictiveCoding/MatrixGaussianChain.lean}{500}{matrixPC\_not\_equilibrium\_of\_ne\_posteriorMean} & (negative fixture) \\
  & \leansource{lean/mettapedia/Mettapedia/MachineLearning/NeuralNetworks/PredictiveCoding/LinearGaussianChain.lean}{1070}{pcConditionalPosteriorMean\_equilibrium} & LinearGaussian\-Chain (scalar) \\
Thm~\ref{thm:errrec}
  & \leansource{lean/mettapedia/Mettapedia/MachineLearning/NeuralNetworks/PredictiveCoding/BackpropExactness.lean}{28}{equilibriumError\_satisfies\_backpropRecursion} & BackpropExactness \\
  & \leansource{lean/mettapedia/Mettapedia/MachineLearning/NeuralNetworks/PredictiveCoding/BackpropExactness.lean}{124}{tiedEquilibriumError\_satisfies\_bpttRecursion} & (weight tying) \\
Thm~\ref{thm:dagrec}
  & \leansource{lean/mettapedia/Mettapedia/MachineLearning/NeuralNetworks/PredictiveCoding/SharedLatentDAG.lean}{223}{sharedDAGEquilibriumError\_satisfies\_parentAggregationRecursion}
  & SharedLatentDAG \\
  & \leansource{lean/mettapedia/Mettapedia/MachineLearning/NeuralNetworks/PredictiveCoding/SharedLatentDAG.lean}{324}{dagMultiParent\_parentAggregate\_counts\_both\_slots} & (positive fixture) \\
  & \leansource{lean/mettapedia/Mettapedia/MachineLearning/NeuralNetworks/PredictiveCoding/SharedLatentDAG.lean}{333}{dagMultiParent\_single\_slot\_is\_incorrect\_negative\_example}
  & (negative fixture) \\
Thm~\ref{thm:gradne}
  & \leansource{lean/mettapedia/Mettapedia/MachineLearning/NeuralNetworks/PredictiveCoding/BackpropExactness.lean}{101}{equilibriumGradient\_ne\_backpropGradient\_example} & BackpropExactness \\
  & \leansource{lean/mettapedia/Mettapedia/MachineLearning/NeuralNetworks/PredictiveCoding/BackpropExactness.lean}{89}{equilibrium\_state\_ne\_forward\_state\_example} & (root cause) \\
Thm~\ref{thm:zil}
  & \leansource{lean/mettapedia/Mettapedia/MachineLearning/NeuralNetworks/PredictiveCoding/ZILExactness.lean}{190}{zilScheduledGainGradient\_eq\_backpropGradient} & ZILExactness \\
  & \leansource{lean/mettapedia/Mettapedia/MachineLearning/NeuralNetworks/PredictiveCoding/ZILExactness.lean}{158}{zilChainResidual\_at\_schedule\_eq\_backpropError} & (front lemma) \\
  & \leansource{lean/mettapedia/Mettapedia/MachineLearning/NeuralNetworks/PredictiveCoding/ZILExactness.lean}{230}{zilOffSchedule\_freeEquilibriumGradient\_failure} & (negative) \\
  & \leansource{lean/mettapedia/Mettapedia/MachineLearning/NeuralNetworks/PredictiveCoding/ZILExactness.lean}{251}{residualSkip\_naiveFirstArrivalGradient\_failure\_example} & (negative) \\
Thm~\ref{thm:signed}
  & \leansource{lean/mettapedia/Mettapedia/MachineLearning/NeuralNetworks/PredictiveCoding/DAGScheduleExactness.lean}{204}{dagScheduleExactAtEveryNode\_iff\_all\_missed\_eq\_zero}
  & DAGSchedule\-Exactness \\
  & \leansource{lean/mettapedia/Mettapedia/MachineLearning/NeuralNetworks/PredictiveCoding/DAGScheduleExactness.lean}{218}{dagScheduleExactAtEveryNode\_iff\_each\_admissible\_or\_cancellation} & \\
  & \leansource{lean/mettapedia/Mettapedia/MachineLearning/NeuralNetworks/PredictiveCoding/DAGScheduleExactness.lean}{233}{dagSchedule\_exact\_not\_admissible\_iff\_hasMissedCancellation} & \\
  & \leansource{lean/mettapedia/Mettapedia/MachineLearning/NeuralNetworks/PredictiveCoding/DAGScheduleExactness.lean}{275}{dagScheduleExactAtEveryNode\_iff\_admissible\_of\_noCancellation}
  & (fusion under no cancellation) \\
Thm~\ref{thm:structsigned}
  & \leansource{lean/mettapedia/Mettapedia/MachineLearning/NeuralNetworks/PredictiveCoding/DAGScheduleExactness.lean}{411}{dagScheduleUniversallyExactAtEveryNode\_iff\_targetGain\_zero}
  & DAGSchedule\-Exactness \\
  & \leansource{lean/mettapedia/Mettapedia/MachineLearning/NeuralNetworks/PredictiveCoding/DAGScheduleExactness.lean}{609}{dagFirstArrivalScheduleUniversallyExact\_iff\_targetGain\_zero} & \\
  & \leansource{lean/mettapedia/Mettapedia/MachineLearning/NeuralNetworks/PredictiveCoding/ModuleScheduleExactness.lean}{134}{moduleScheduleUniversallyExactAtNode\_iff\_targetGain\_zero}
  & ModuleSchedule\-Exactness \\
  & \leansource{lean/mettapedia/Mettapedia/MachineLearning/NeuralNetworks/PredictiveCoding/ModuleScheduleExactness.lean}{167}{moduleSignedCancellation\_universallyExact\_positive\_example}
  & (positive fixture) \\
  & \leansource{lean/mettapedia/Mettapedia/MachineLearning/NeuralNetworks/PredictiveCoding/ModuleScheduleExactness.lean}{179}{moduleSingleMissedEdge\_not\_universallyExact\_negative\_example}
  & (negative fixture) \\
Thm~\ref{thm:sched}
  & \leansource{lean/mettapedia/Mettapedia/MachineLearning/NeuralNetworks/PredictiveCoding/DAGScheduleExactness.lean}{727}{dagFirstArrivalScheduleExact\_iff\_uniformOutgoingTargetRanks}
  & DAGSchedule\-Exactness \\
  & \leansource{lean/mettapedia/Mettapedia/MachineLearning/NeuralNetworks/PredictiveCoding/DAGScheduleExactness.lean}{873}{uniformLongSkipGraph\_firstArrival\_exact} & (witness) \\
  & \leansource{lean/mettapedia/Mettapedia/MachineLearning/NeuralNetworks/PredictiveCoding/DAGScheduleExactness.lean}{849}{uniformLongSkipGraph\_not\_unitLevelized} & (witness) \\
Thm~\ref{thm:epc}
  & \leansource{lean/mettapedia/Mettapedia/MachineLearning/NeuralNetworks/PredictiveCoding/ErrorStateReparameterization.lean}{40}{spcErrorStream\_epcStateStream} & ErrorState\-Reparameterization \\
  & \leansource{lean/mettapedia/Mettapedia/MachineLearning/NeuralNetworks/PredictiveCoding/ErrorStateReparameterization.lean}{48}{epcStateStream\_spcErrorStream} & \\
  & \leansource{lean/mettapedia/Mettapedia/MachineLearning/NeuralNetworks/PredictiveCoding/ErrorStateReparameterization.lean}{82}{epcPrefixEnergy\_eq\_spcPrefixEnergy} & \\
  & \leansource{lean/mettapedia/Mettapedia/MachineLearning/NeuralNetworks/PredictiveCoding/ErrorStateReparameterization.lean}{138}{epcCriticalDifferential\_iff\_spcCriticalDifferential} & \\
  & \leansource{lean/mettapedia/Mettapedia/MachineLearning/NeuralNetworks/PredictiveCoding/ErrorStateReparameterization.lean}{240}{epcLocalParameterUpdate\_eq\_spcLocalParameterUpdate} & \\
Thm~\ref{thm:wave}
  & \leansource{lean/mettapedia/Mettapedia/MachineLearning/NeuralNetworks/PredictiveCoding/ErrorStateReparameterization.lean}{171}{spcFirstArrivalSignal\_layer\_closedForm}
  & ErrorState\-Reparameterization \\
  & \leansource{lean/mettapedia/Mettapedia/MachineLearning/NeuralNetworks/PredictiveCoding/ErrorStateReparameterization.lean}{185}{spcWavefront\_eventually\_below\_tolerance} & \\
  & \leansource{lean/mettapedia/Mettapedia/MachineLearning/NeuralNetworks/PredictiveCoding/ErrorStateReparameterization.lean}{209}{spcWavefront\_nonzero\_but\_below\_tolerance\_example} & (positive) \\
  & \leansource{lean/mettapedia/Mettapedia/MachineLearning/NeuralNetworks/PredictiveCoding/ErrorStateReparameterization.lean}{218}{spcWavefront\_unitRate\_not\_below\_small\_tolerance} & (boundary) \\
  & \leansource{lean/mettapedia/Mettapedia/MachineLearning/NeuralNetworks/PredictiveCoding/ErrorStateReparameterization.lean}{260}{epcOneStepParameterUpdate\_unitRate\_eq\_backprop} & \\
Thm~\ref{thm:mupcdepth}
  & \leansource{lean/mettapedia/Mettapedia/MachineLearning/NeuralNetworks/PredictiveCoding/MuPCDepthCredit.lean}{54}{spcFirstArrivalParameterUpdate\_eq\_depthCreditBackprop}
  & MuPCDepthCredit \\
  & \leansource{lean/mettapedia/Mettapedia/MachineLearning/NeuralNetworks/PredictiveCoding/MuPCDepthCredit.lean}{77}{chainLinkLoss\_after\_uniformDepthCredit\_exact} & \\
  & \leansource{lean/mettapedia/Mettapedia/MachineLearning/NeuralNetworks/PredictiveCoding/MuPCDepthCredit.lean}{255}{inferenceRateEqualization\_forces\_unitBackpropBoundary} & \\
  & \leansource{lean/mettapedia/Mettapedia/MachineLearning/NeuralNetworks/PredictiveCoding/MuPCDepthCredit.lean}{366}{muPCDepthScaling\_cannot\_cancel\_contractingDepthCredit} & \\
  & \leansource{lean/mettapedia/Mettapedia/MachineLearning/NeuralNetworks/PredictiveCoding/MuPCDepthCredit.lean}{604}{uniformPrecision\_equilibrium\_sub\_forward\_exact} & \\
Thm~\ref{thm:depth}
  & \leansource{lean/mettapedia/Mettapedia/MachineLearning/NeuralNetworks/PredictiveCoding/DepthPathology.lean}{85}{pathSlowMode\_eigenrelation} & DepthPathology \\
  & \leansource{lean/mettapedia/Mettapedia/MachineLearning/NeuralNetworks/PredictiveCoding/DepthPathology.lean}{106}{pathSlowMode\_iterate\_exact} & \\
  & \leansource{lean/mettapedia/Mettapedia/MachineLearning/NeuralNetworks/PredictiveCoding/DepthPathology.lean}{139}{pathRelaxationRate\_gap\_bound} & \\
  & \leansource{lean/mettapedia/Mettapedia/MachineLearning/NeuralNetworks/PredictiveCoding/DepthPathology.lean}{152}{pathRelaxationRate\_tendsto\_one} & \\
Thm~\ref{thm:het}
  & \leansource{lean/mettapedia/Mettapedia/MachineLearning/NeuralNetworks/PredictiveCoding/HeterogeneousDepthBlocks.lean}{114}{HeterogeneousDepthBlockFamily.modalSpectralRadius\_lt\_one\_iff}
  & Heterogeneous\-DepthBlocks \\
  & \leansource{lean/mettapedia/Mettapedia/MachineLearning/NeuralNetworks/PredictiveCoding/HeterogeneousDepthBlocks.lean}{152}{HeterogeneousDepthBlockFamily.modalSpectralGap\_le\_depth\_bound} & \\
Thm~\ref{cor:adapter}
  & \leansource{lean/mettapedia/Mettapedia/MachineLearning/NeuralNetworks/PredictiveCoding/FrozenAdapterRestriction.lean}{67}{coupledBackboneAdapterStep\_not\_blockDiagonal}
  & FrozenAdapter\-Restriction \\
  & \leansource{lean/mettapedia/Mettapedia/MachineLearning/NeuralNetworks/PredictiveCoding/FrozenAdapterRestriction.lean}{152}{frozenAdapterRestriction\_eq\_blockPathJacobiStep\_clamped} & \\
  & \leansource{lean/mettapedia/Mettapedia/MachineLearning/NeuralNetworks/PredictiveCoding/FrozenAdapterRestriction.lean}{189}{frozenAdapterRestriction\_eq\_blockPathJacobiStep\_of\_boundary\_clamped} & \\
  & \leansource{lean/mettapedia/Mettapedia/MachineLearning/NeuralNetworks/PredictiveCoding/FrozenAdapterRestriction.lean}{233}{isFrozenAdapterSlowModeRate\_iff} & \\
  & \leansource{lean/mettapedia/Mettapedia/MachineLearning/NeuralNetworks/PredictiveCoding/FrozenAdapterRestriction.lean}{281}{frozenAdapterRestriction\_gap\_bound} & \\
  & \leansource{lean/mettapedia/Mettapedia/MachineLearning/NeuralNetworks/PredictiveCoding/FrozenAdapterRestriction.lean}{336}{activeBackboneAdapter\_backbone\_gap\_bound} & \\
  & \leansource{lean/mettapedia/Mettapedia/MachineLearning/NeuralNetworks/PredictiveCoding/FrozenAdapterRestriction.lean}{394}{partialFreezing\_reimports\_backbone\_negative\_fixture} & \\
Thm~\ref{thm:precond}
  & \leansource{lean/mettapedia/Mettapedia/MachineLearning/NeuralNetworks/PredictiveCoding/PreconditionedLearning.lean}{99}{chainLinkLoss\_preconditionedUpdate\_exact} & Preconditioned\-Learning \\
  & \leansource{lean/mettapedia/Mettapedia/MachineLearning/NeuralNetworks/PredictiveCoding/PreconditionedLearning.lean}{81}{chainLinkResidual\_preconditionedUpdate\_exact} & \\
  & \leansource{lean/mettapedia/Mettapedia/MachineLearning/NeuralNetworks/PredictiveCoding/PreconditionedLearning.lean}{196}{chainLink\_zeroLoss\_afterUpdate\_iff\_identity} & \\
  & \leansource{lean/mettapedia/Mettapedia/MachineLearning/NeuralNetworks/PredictiveCoding/PreconditionedLearning.lean}{166}{chainLink\_preconditionedUpdate\_strictDescent\_iff} & \\
  & \leansource{lean/mettapedia/Mettapedia/MachineLearning/NeuralNetworks/PredictiveCoding/PreconditionedLearning.lean}{123}{chainLink\_lossDecrease\_gap\_from\_identity\_exact} & \\
  & \leansource{lean/mettapedia/Mettapedia/MachineLearning/NeuralNetworks/PredictiveCoding/PreconditionedLearning.lean}{148}{chainLink\_lossDescent\_strictly\_degrades\_with\_squaredDeparture} & \\
  & \leansource{lean/mettapedia/Mettapedia/MachineLearning/NeuralNetworks/PredictiveCoding/PreconditionedLearning.lean}{242}{unitChain\_large\_preconditioner\_increases\_loss\_negative\_example} & \\
Thm~\ref{thm:spectral}
  & \leansource{lean/mettapedia/Mettapedia/MachineLearning/NeuralNetworks/PredictiveCoding/OperatorStability.lean}{22}{spectralRadius\_lt\_one\_iff\_pow\_tendsto\_zero} & Operator\-Stability \\
  & \leansource{lean/mettapedia/Mettapedia/MachineLearning/NeuralNetworks/PredictiveCoding/OperatorStability.lean}{233}{complexMatrixPreconditioner\_fullSpectrum\_convergence} & \\
  & \leansource{lean/mettapedia/Mettapedia/MachineLearning/NeuralNetworks/PredictiveCoding/OperatorStability.lean}{416}{nonnormalTransientMatrix\_fullSpectrum\_stable} & (nonnormal witness) \\
  & \leansource{lean/mettapedia/Mettapedia/MachineLearning/NeuralNetworks/PredictiveCoding/OperatorStability.lean}{404}{nonnormalTransientMatrix\_not\_normal} & \\
Thm~\ref{thm:envelope}
  & \leansource{lean/mettapedia/Mettapedia/MachineLearning/NeuralNetworks/PredictiveCoding/OperatorStability.lean}{79}{operatorPowNorm\_geometricEnvelope} & Operator\-Stability \\
  & \leansource{lean/mettapedia/Mettapedia/MachineLearning/NeuralNetworks/PredictiveCoding/OperatorStability.lean}{141}{operatorResidualIterate\_norm\_geometricEnvelope} & \\
Thm~\ref{thm:jordan}
  & \leansource{lean/mettapedia/Mettapedia/MachineLearning/NeuralNetworks/PredictiveCoding/OperatorStability.lean}{255}{jordanTransientMatrix\_pow\_exact} & Operator\-Stability \\
  & \leansource{lean/mettapedia/Mettapedia/MachineLearning/NeuralNetworks/PredictiveCoding/OperatorStability.lean}{280}{jordanTransientMatrix\_spectrum} & \\
  & \leansource{lean/mettapedia/Mettapedia/MachineLearning/NeuralNetworks/PredictiveCoding/OperatorStability.lean}{304}{jordanTransientMatrix\_not\_normal} & \\
  & \leansource{lean/mettapedia/Mettapedia/MachineLearning/NeuralNetworks/PredictiveCoding/OperatorStability.lean}{364}{spectralRadius\_zero\_allows\_arbitrary\_transientAmplification}
  & (negative boundary) \\
Thm~\ref{thm:mode}
  & \leansource{lean/mettapedia/Mettapedia/MachineLearning/NeuralNetworks/PredictiveCoding/PreconditionedLearning.lean}{380}{matrixPreconditionedResidual\_loss\_perMode\_exact}
  & Preconditioned\-Learning \\
  & \leansource{lean/mettapedia/Mettapedia/MachineLearning/NeuralNetworks/PredictiveCoding/PreconditionedLearning.lean}{542}{mixedModePreconditioner\_helps\_one\_hurts\_another\_fixture} & \\
Thm~\ref{thm:armijo}
  & \leansource{lean/mettapedia/Mettapedia/MachineLearning/NeuralNetworks/PredictiveCoding/BacktrackingDescent.lean}{65}{armijoBacktracking\_terminates} & Backtracking\-Descent \\
  & \leansource{lean/mettapedia/Mettapedia/MachineLearning/NeuralNetworks/PredictiveCoding/BacktrackingDescent.lean}{78}{armijoCondition\_energy\_nonincreasing} & \\
  & \leansource{lean/mettapedia/Mettapedia/MachineLearning/NeuralNetworks/PredictiveCoding/BacktrackingDescent.lean}{150}{pcEnergyDepth2\_armijoBacktracking\_terminates\_with\_descent} & \\
Thm~\ref{thm:saddles}
  & \leansource{lean/mettapedia/Mettapedia/MachineLearning/NeuralNetworks/PredictiveCoding/EquilibratedEnergySaddles.lean}{167}{orthogonalDataset\_outputDirection\_curvature}
  & EquilibratedEnergy\-Saddles \\
  & \leansource{lean/mettapedia/Mettapedia/MachineLearning/NeuralNetworks/PredictiveCoding/EquilibratedEnergySaddles.lean}{236}{orthogonalDataset\_innerStrip\_hasStrictSaddleCertificate} & \\
  & \leansource{lean/mettapedia/Mettapedia/MachineLearning/NeuralNetworks/PredictiveCoding/EquilibratedEnergySaddles.lean}{247}{orthogonalDataset\_counterexamplePoint\_all\_curvatures\_zero} & \\
  & \leansource{lean/mettapedia/Mettapedia/MachineLearning/NeuralNetworks/PredictiveCoding/EquilibratedEnergySaddles.lean}{282}{orthogonalDataset\_counterexamplePoint\_twoSidedSaddleCurves} & \\
  & \leansource{lean/mettapedia/Mettapedia/MachineLearning/NeuralNetworks/PredictiveCoding/EquilibratedEnergySaddles.lean}{319}{orthogonalDataset\_universalStrictSaddleConjecture\_counterexample}
  & (negative boundary) \\
Thm~\ref{thm:coupled}
  & \leansource{lean/mettapedia/Mettapedia/MachineLearning/NeuralNetworks/PredictiveCoding/CoupledDynamics.lean}{198}{coupledTwoTimescale\_multipliers\_stable\_iff} & CoupledDynamics \\
  & \leansource{lean/mettapedia/Mettapedia/MachineLearning/NeuralNetworks/PredictiveCoding/CoupledDynamics.lean}{288}{coupledTwoTimescale\_converges\_below\_threshold} & \\
  & \leansource{lean/mettapedia/Mettapedia/MachineLearning/NeuralNetworks/PredictiveCoding/CoupledDynamics.lean}{325}{coupledTwoTimescale\_diverges\_above\_learning\_threshold} & \\
Thm~\ref{thm:turnover}
  & \leansource{lean/mettapedia/Mettapedia/MachineLearning/NeuralNetworks/PredictiveCoding/EntropyTurnover.lean}{280}{arbitraryDepthLinearEntropyTurnover\_general} & EntropyTurnover \\
  & \leansource{lean/mettapedia/Mettapedia/MachineLearning/NeuralNetworks/PredictiveCoding/EntropyTurnover.lean}{232}{arbitraryDepthTurnoverEntropy\_nonmonotone\_iff} & \\
Thm~\ref{thm:nonlin}
  & \leansource{lean/mettapedia/Mettapedia/MachineLearning/NeuralNetworks/PredictiveCoding/EntropyTurnover.lean}{630}{monotoneContractiveNonlinearEntropyTurnover\_sufficient} & EntropyTurnover \\
  & \leansource{lean/mettapedia/Mettapedia/MachineLearning/NeuralNetworks/PredictiveCoding/EntropyTurnover.lean}{550}{MonotoneContractiveSettling.iterate\_abs\_le} & (contraction bound) \\
  & \leansource{lean/mettapedia/Mettapedia/MachineLearning/NeuralNetworks/PredictiveCoding/EntropyTurnover.lean}{678}{monotoneContractiveNonlinearEntropyTurnover\_uniform\_boundary} & (boundary) \\
Thm~\ref{thm:trainednonlin}
  & \leansource{lean/mettapedia/Mettapedia/MachineLearning/NeuralNetworks/PredictiveCoding/TrainedNonlinearTurnover.lean}{208}{trainedCenteredSigmoidWeight\_is\_unique\_global\_minimizer}
  & TrainedNonlinear\-Turnover \\
  & \leansource{lean/mettapedia/Mettapedia/MachineLearning/NeuralNetworks/PredictiveCoding/TrainedNonlinearTurnover.lean}{241}{trainedCenteredSigmoidEntropyTurnover} & \\
  & \leansource{lean/mettapedia/Mettapedia/MachineLearning/NeuralNetworks/PredictiveCoding/TrainedNonlinearTurnover.lean}{312}{approximateTraining\_preservesCenteredSigmoidTurnoverMargin} & \\
  & \leansource{lean/mettapedia/Mettapedia/MachineLearning/NeuralNetworks/PredictiveCoding/TrainedNonlinearTurnover.lean}{270}{centeredSigmoidContraction\_does\_not\_reach\_logit\_one} & (boundary) \\
  & \leansource{lean/mettapedia/Mettapedia/MachineLearning/NeuralNetworks/PredictiveCoding/TrainedNonlinearTurnover.lean}{395}{trainedCenteredSigmoidVectorTurnover} & (vector lift) \\
  & \leansource{lean/mettapedia/Mettapedia/MachineLearning/NeuralNetworks/PredictiveCoding/TrainedNonlinearTurnover.lean}{416}{centeredSigmoidVectorZeroSeed\_uniform} & (negative fixture) \\
Appendix
  & \leansource{lean/mettapedia/Mettapedia/MachineLearning/NeuralNetworks/PredictiveCoding/CausalDirectionIdentifiability.lean}{210}{PositiveBivariateCovariance.\allowbreak distinct\_orientations\_same\_observational\_batchEnergy}
  & CausalDirection\-Identifiability \\
  & \leansource{lean/mettapedia/Mettapedia/MachineLearning/NeuralNetworks/PredictiveCoding/CausalDirectionIdentifiability.lean}{683}{PositiveBivariateCovariance.\allowbreak clamp\_is\_conditioning\_and\_direction\_blind} & \\
  & \leansource{lean/mettapedia/Mettapedia/MachineLearning/NeuralNetworks/PredictiveCoding/CausalDirectionIdentifiability.lean}{759}{PositiveBivariateCovariance.\allowbreak interventionalEquilibria\_separate\_iff} & \\
  & \leansource{lean/mettapedia/Mettapedia/MachineLearning/NeuralNetworks/PredictiveCoding/CausalDirectionIdentifiability.lean}{866}{PositiveBivariateCovariance.\allowbreak observationally\_blind\_intervention\_certificate} & \\
  & \leansource{lean/mettapedia/Mettapedia/MachineLearning/NeuralNetworks/PredictiveCoding/CausalDirectionIdentifiability.lean}{976}{PositiveBivariateCovariance.\allowbreak TwoNodeQuadraticResidualModel.\allowbreak batchEnergy\_eq\_trace\_quadraticForm\_mul\_empiricalSecondMoment} & \\
Thm~\ref{thm:intervention}
  & \leansource{lean/mettapedia/Mettapedia/MachineLearning/NeuralNetworks/PredictiveCoding/CausalDirectionIdentifiability.lean}{853}{PositiveBivariateCovariance.\allowbreak separatesDirectionsByInterventionalEnergy\_iff}
  & CausalDirection\-Identifiability \\
  & \leansource{lean/mettapedia/Mettapedia/MachineLearning/NeuralNetworks/PredictiveCoding/CausalDirectionIdentifiability.lean}{907}{PositiveBivariateCovariance.\allowbreak unitNoiseForwardCovariance\_observation\_tied\_intervention\_separates}
  & (positive fixture) \\
  & \leansource{lean/mettapedia/Mettapedia/MachineLearning/NeuralNetworks/PredictiveCoding/CausalDirectionIdentifiability.lean}{921}{PositiveBivariateCovariance.\allowbreak independentEqualVarianceCovariance\_interventionalEnergy\_fixture}
  & (negative fixture) \\
Thm~\ref{thm:causalselection}
  & \leansource{lean/mettapedia/Mettapedia/MachineLearning/NeuralNetworks/PredictiveCoding/CausalDirectionSelection.lean}{316}{PositiveBivariateCovariance.\allowbreak forwardOrientationWinsByMinimizedInterventionalEnergy\_iff}
  & CausalDirection\-Selection \\
  & \leansource{lean/mettapedia/Mettapedia/MachineLearning/NeuralNetworks/PredictiveCoding/CausalDirectionSelection.lean}{222}{PositiveBivariateCovariance.\allowbreak isMinimizedForwardInterventionalBatchEnergy\_iff}
  & (unique forward fit) \\
  & \leansource{lean/mettapedia/Mettapedia/MachineLearning/NeuralNetworks/PredictiveCoding/CausalDirectionSelection.lean}{282}{PositiveBivariateCovariance.\allowbreak reverseCandidates\_distinct\_but\_intervention\_deletes\_parameter}
  & (anti-vacuity) \\
  & \leansource{lean/mettapedia/Mettapedia/MachineLearning/NeuralNetworks/PredictiveCoding/CausalDirectionSelection.lean}{365}{PositiveBivariateCovariance.\allowbreak unitNoiseForwardCovariance\_minimizedIntervention\_selects\_forward}
  & (positive fixture) \\
  & \leansource{lean/mettapedia/Mettapedia/MachineLearning/NeuralNetworks/PredictiveCoding/CausalDirectionSelection.lean}{375}{PositiveBivariateCovariance.\allowbreak independentEqualVariance\_minimizedIntervention\_tied}
  & (negative fixture) \\
Appendix
  & \leansource{lean/mettapedia/Mettapedia/MachineLearning/ContinualLearning/QuadraticTwoTask.lean}{46}{sequentialTwoTaskUpdate\_sub\_additive\_exact}
  & ContinualLearning/ QuadraticTwoTask \\
  & \leansource{lean/mettapedia/Mettapedia/MachineLearning/ContinualLearning/QuadraticTwoTask.lean}{80}{centered\_sequential\_order\_defect\_eq\_commutator} & \\
  & \leansource{lean/mettapedia/Mettapedia/MachineLearning/ContinualLearning/QuadraticTwoTask.lean}{214}{linearTwoTask\_causalCoding\_separation\_crown} & \\
  & \leansource{lean/mettapedia/Mettapedia/MachineLearning/ContinualLearning/EvidenceLedger.lean}{74}{GaussianEvidence.reused\_contribution\_exact\_excess}
  & ContinualLearning/ EvidenceLedger \\
  & \leansource{lean/mettapedia/Mettapedia/MachineLearning/ContinualLearning/EvidenceLedger.lean}{156}{linearCausalCodingBridge\_crown} & \\
  & \leansource{lean/mettapedia/Mettapedia/MachineLearning/ContinualLearning/EvidenceLedger.lean}{149}{generalCausalCodingCorrespondence\_requires\_interventionalClamp}
  & (scope marker) \\
Thm~\ref{thm:yield}
  & \leansource{lean/mettapedia/Mettapedia/MachineLearning/SearchGuidance/DecisivenessYield.lean}{88}{iidExpectedTargetHits\_eq\_budget\_mul\_targetMass}
  & SearchGuidance/ \mbox{DecisivenessYield} \\
  & \leansource{lean/mettapedia/Mettapedia/MachineLearning/SearchGuidance/DecisivenessYield.lean}{176}{iidExpectedTargetHits\_strictMono\_targetMass} & \\
Thm~\ref{thm:sep}
  & \leansource{lean/mettapedia/Mettapedia/MachineLearning/SearchGuidance/DecisivenessYield.lean}{242}{equal\_crossEntropy\_strictly\_ordered\_expectedYield}
  & SearchGuidance/ \mbox{DecisivenessYield} \\
\bottomrule
\end{longtable}
}

\section*{Related work}

Equivalence-in-the-limit results and their scope: \cite{whittington2017} (MLPs),
\cite{millidge2020} (arbitrary graphs, including LSTMs at small scale), \cite{song2020}
(exact scheduling), \cite{salvatori2022} (topological limits). Engineering of deep PC:
\cite{goemaere2025}, whose diagnosis that naive settling misbehaves in digital simulation
Theorems~\ref{thm:depth} and~\ref{thm:precond} corroborate from the analytic side.
Equilibrated-energy landscape results and the universal strict-saddle conjecture are due
to \cite{innocenti2025}; Section~\ref{sec:saddles} proves a degenerate boundary outside
their covered class. The
program-synthesis substrate is Gauthier's OEIS self-learning system and its NMT variant
\cite{gauthier2023}. Our companion report \emph{AlienCoding} documents the replication
methodology, the full campaign, and the reference numbers.

\begin{thebibliography}{9}

\bibitem{whittington2017} J.~Whittington and R.~Bogacz.
An approximation of the error backpropagation algorithm in a predictive coding network
with local Hebbian synaptic plasticity. \emph{Neural Computation}, 29(5), 2017.

\bibitem{song2020} Y.~Song, T.~Lukasiewicz, Z.~Xu, and R.~Bogacz.
Can the brain do backpropagation? Exact implementation of backpropagation in predictive
coding networks. \emph{NeurIPS}, 2020.

\bibitem{millidge2020} B.~Millidge, A.~Tschantz, and C.~Buckley.
Predictive coding approximates backprop along arbitrary computation graphs.
arXiv:2006.04182, 2020.

\bibitem{salvatori2022} T.~Salvatori et al.
Learning on arbitrary graph topologies via predictive coding. \emph{NeurIPS}, 2022.

\bibitem{goemaere2025} C.~Goemaere, G.~Oliviers, R.~Bogacz, and T.~Demeester.
ePC: Fast and deep predictive coding in digital simulation. arXiv:2505.20137, 2025.

\bibitem{innocenti2025} F.~Innocenti, E.~M.~Achour, R.~Singh, and C.~L.~Buckley.
Only strict saddles in the energy landscape of predictive coding networks?
\emph{Journal of Statistical Mechanics: Theory and Experiment}, 2025.
doi:10.1088/1742-5468/ade2eb.

\bibitem{innocentimupc2025} F.~Innocenti, E.~M.~Achour, and C.~L.~Buckley.
$\mu$PC: Scaling predictive coding to 100+ layer networks.
arXiv:2505.13124v2, 2025.

\bibitem{pearl2012} J.~Pearl.
The do-calculus revisited.
In \emph{Proceedings of the Twenty-Eighth Conference on Uncertainty in Artificial
Intelligence}, 2012.

\bibitem{halpern2000} J.~Y.~Halpern.
Axiomatizing causal reasoning.
\emph{Journal of Artificial Intelligence Research}, 12:317--337, 2000.

\bibitem{gauthier2023} T.~Gauthier, M.~Ol\v{s}\'ak, and J.~Urban.
Alien coding. arXiv:2301.11479, 2023.

\end{thebibliography}

\end{document}
